\documentclass[11pt,a4paper]{article}

% ============================================================
% Packages (core only --- texlive-latex-base + texlive-base)
% ============================================================
\usepackage{amsmath, amssymb, amsthm}
\usepackage[a4paper, margin=2.5cm]{geometry}
\usepackage{array}

% ============================================================
% Theorem environments
% ============================================================
\theoremstyle{definition}
\newtheorem{definition}{Definition}[section]
\newtheorem{remark}{Remark}[section]
\newtheorem{proposition}{Proposition}[section]

\theoremstyle{plain}
\newtheorem{theorem}{Theorem}[section]
\newtheorem{lemma}{Lemma}[section]

% ============================================================
% Notation macros
% ============================================================
\newcommand{\R}{\mathbb{R}}
\newcommand{\E}{\mathbb{E}}
\newcommand{\N}{\mathbb{N}}
\newcommand{\norm}[1]{\left\lVert #1 \right\rVert}
\newcommand{\inner}[2]{\langle #1,\, #2 \rangle}
\newcommand{\corr}{\operatorname{corr}}
\newcommand{\med}{\operatorname{median}}
\newcommand{\sgn}{\operatorname{sign}}
\DeclareMathOperator{\IC}{IC}
\DeclareMathOperator*{\argmax}{arg\,max}
\DeclareMathOperator*{\argmin}{arg\,min}
\newcommand{\code}[1]{\texttt{\detokenize{#1}}}

% ============================================================
\begin{document}

\title{\textbf{Autonomous Microstructure Alpha Discovery\\[3pt]
under False-Discovery Control}\\[6pt]
\large Order-Book Imbalance, Conditional Information Coefficients,\\
and the Adverse-Selection Barrier in Hyperliquid Perpetual Futures}

\author{Yigit Onat\\
\textit{Independent Researcher}\\
\texttt{yionat@gmail.com}}
\date{June 2026}
\maketitle

% ============================================================
\begin{abstract}
We study the gap between statistical predictability and economic exploitability in a
decentralized derivatives venue. Using $100$\,ms order-book and trade data for three
Hyperliquid perpetual futures (BTC, ETH, SOL), we compute the Spearman information
coefficient (IC) of $236$ engineered microstructure features against signed and absolute
forward returns at horizons from one second to fifteen minutes. Discovery and validation are
performed by an autonomous platform that subjects every candidate to a five-gate replication
protocol (discovery, cost, temporal, cross-symbol, and correlation gates) with
Benjamini--Hochberg false-discovery-rate control, an adaptive acceptance threshold, and
continuous decay monitoring---a disciplined apparatus designed to suppress the multiple-testing
bias that pervades empirical asset pricing. Order-book imbalance is the dominant directional
predictor, attaining $\IC\approx 0.45$ at the $1$--$5$\,second horizon and decaying to near zero
by five minutes, with the same feature hierarchy across all three symbols and a clean
directional/volatility orthogonality. The signal passes seven of eight robustness checks for
BTC and six of eight for ETH and SOL. Our central result is a negative one: the
\emph{conditional} information coefficient, measured on the subset of states where a
directionally-correct passive (mid-cross) fill would occur, collapses from $0.45$ to
$\approx 0.03$ across all symbols. This is a structural adverse-selection barrier, not an
estimation artifact, and it identifies execution---not signal quality---as the binding
constraint. We then present a three-layer hierarchical combiner that gates fast microstructure
timing on a slow directional bias, recovering positive cost-adjusted out-of-sample performance
($\IC$ of $0.18$/$0.25$/$0.36$ for BTC/ETH/SOL) on a thin window. The contribution is twofold: a
reproducible methodology for false-discovery-controlled microstructure research, and direct
evidence that a strong, universal, robust predictor can nonetheless be economically inert under
realistic passive execution.
\end{abstract}

\smallskip
\noindent\textbf{Keywords:} market microstructure, information coefficient, order-book imbalance,
adverse selection, false discovery rate, cryptocurrency, perpetual futures, automated research

\smallskip
\noindent\textbf{JEL Classification:} C12, C55, G14, G17

\tableofcontents
\bigskip

% ============================================================
\section{Introduction}
\label{sec:intro}
% ============================================================

\subsection{Predictability versus profitability}

A recurring failure mode in empirical finance is to mistake a statistically significant
predictor for a tradeable one. A feature may correlate strongly with future returns in sample,
survive out-of-sample testing, and still yield no economic profit once the mechanics of
execution are accounted for. This paper provides a clean, quantified instance of that gap in a
high-frequency setting, and argues that the gap---not the discovery of predictors---is where the
binding constraint lies.

We study Hyperliquid perpetual futures, a decentralized derivatives exchange, at $100$\,ms
resolution. We find that order-book imbalance predicts the direction of returns with an
information coefficient (IC) of approximately $0.45$ at the $1$--$5$ second horizon---a very
strong microstructure signal that is universal across BTC, ETH, and SOL and robust across days,
intraday windows, and volatility regimes. Yet when we condition the IC on the event that a
directionally-correct passive limit order would actually fill, the coefficient collapses to
$\approx 0.03$. The predictability is real; the profitability, under a passive mid-cross fill
model, is not.

\subsection{The horizon mistake}

The investigation began from a diagnostic failure: a library of roughly twenty-five trading
algorithms, all evaluated at a $100$-minute forecast horizon using five-minute bar aggregation,
were uniformly unprofitable. The cause was an aggregation artifact. The predictive content of
microstructure features lives at the $1$--$5$ second scale; aggregating to five-minute bars and
forecasting one hundred minutes ahead destroys $\IC\approx 0.45$ down to a weak residual
($\IC<0.07$). The features were never the problem---the evaluation horizon was. This reframing
motivates a systematic, horizon-resolved IC study of the full feature set.

\subsection{An autonomous discovery apparatus}

Microstructure research at this scale is acutely vulnerable to data-snooping: with hundreds of
features, multiple horizons, three symbols, and many candidate transformations, the number of
implicit hypotheses is large, and ordinary $p$-values are not credible. We therefore conduct
discovery through an autonomous platform whose design is itself a contribution. Every candidate
signal must independently clear five replication gates---discovery (IC and its incremental
component), cost, temporal replication, cross-symbol replication, and correlation
deduplication---after which Benjamini--Hochberg false-discovery-rate (FDR) control is applied per
cycle, the acceptance threshold adapts upward as strong signals accumulate, and registered
signals are continuously monitored for decay. The platform ingests order-book and trade data in
Rust at $100$\,ms and runs the discovery loop without human intervention. We treat it here as the
\emph{experimental apparatus}: the empirical results below are the output of a discovery process
explicitly engineered to make false discoveries unlikely.

\subsection{Contributions}

\begin{enumerate}
  \item A horizon-resolved information-coefficient map of $236$ microstructure features across
        three cryptocurrency perpetuals, establishing order-book imbalance as the dominant
        directional predictor ($\IC\approx 0.45$ at $1$--$5$\,s, half-life $\approx 15$--$30$\,s),
        with eight approximately independent directional axes and a clean
        directional/volatility orthogonality.
  \item A battery of eight robustness checks (per-day, intraday, volatility-regime, bootstrap,
        temporal out-of-sample) demonstrating that the signal is real, stable, and universal.
  \item The central finding: a \emph{conditional} information coefficient that isolates the
        adverse-selection barrier. Conditioning on directionally-correct passive fills collapses
        $\IC$ from $0.45$ to $\approx 0.03$, identifying execution as the binding constraint.
  \item A three-layer hierarchical combiner that gates fast timing on a slow directional bias and
        recovers positive cost-adjusted out-of-sample performance, the first architecture in this
        programme to confront adverse selection directly.
  \item A reproducible methodology for false-discovery-controlled microstructure research---a
        five-gate replication protocol with FDR control, adaptive thresholds, and
        information-theoretic feature selection.
\end{enumerate}

% ============================================================
\section{Related Work}
\label{sec:related}
% ============================================================

\paragraph{Order-book imbalance and price prediction.}
The predictive content of limit-order-book imbalance is well established. Cont, Stoikov, and
Talreja (2010) model the order book as a continuous-time Markov process and link imbalance to
short-horizon price moves; Cont, Kukanov, and Stoikov (2014) formalize order-flow imbalance as a
linear driver of price changes. Our contribution is not the existence of the imbalance signal but
its precise horizon profile, cross-symbol universality, and---decisively---its behavior
conditional on execution.

\paragraph{Adverse selection and the cost of liquidity provision.}
The classical microstructure literature explains why a passive provider cannot trivially harvest
a directional signal. Glosten and Milgrom (1985) and Kyle (1985) show that informed order flow
imposes adverse selection on liquidity providers; a quote that fills does so disproportionately
when the provider is on the wrong side. Our conditional-IC result is a direct, model-free
measurement of this effect: the very imbalance that predicts an up-move is the imbalance under
which a passive buy fills only after the move has reversed.

\paragraph{Flow toxicity.}
Easley, L\'opez de Prado, and O'Hara (2012) introduce VPIN as a measure of order-flow toxicity in
high-frequency markets. Toxicity is the aggregate counterpart of the per-fill adverse selection
we isolate; we find, consistently, that imbalance carries directional information while toxicity
and intensity features carry volatility (magnitude) information, with near-zero cross-loading.

\paragraph{Multiple testing in empirical finance.}
Harvey, Liu, and Zhu (2016) document rampant multiple testing in the cross-section of returns and
call for higher significance hurdles; Bailey and L\'opez de Prado (2014) introduce the deflated
Sharpe ratio to correct for selection bias and backtest overfitting; White (2000) provides the
reality-check bootstrap. Our discovery apparatus operationalizes this agenda: Benjamini and
Hochberg (1995) FDR control, an adaptive acceptance threshold, and mandatory temporal and
cross-symbol replication are built into the pipeline rather than applied post hoc.

\paragraph{Information-theoretic dependence.}
Kraskov, St\"ogbauer, and Grassberger (2004) provide the $k$-nearest-neighbor mutual-information
estimator that underlies our information-theoretic feature selection; Schreiber (2000) introduces
transfer entropy for directed dependence. We use these as a complement to rank correlation, with
a cost-aware minimum-information threshold derived below.

\paragraph{Cryptocurrency microstructure.}
Makarov and Schoar (2020) study price formation and arbitrage across crypto venues; A\"it-Sahalia
and Brunetti (2024) analyze order-flow toxicity in Bitcoin. We contribute a systematic,
false-discovery-controlled microstructure study on a decentralized perpetual-futures venue, a
market structure distinct from the centralized spot exchanges that dominate the literature.

% ============================================================
\section{Data and the Discovery Apparatus}
\label{sec:data}
% ============================================================

\subsection{Market data}

The data are Hyperliquid perpetual futures for three symbols (BTC, ETH, SOL) sampled at
$100$\,ms. A Rust ingestor subscribes to the Level-2 WebSocket feed and computes a
$236$-dimensional feature vector per symbol every $100$\,ms, written to columnar storage. The
primary discovery window is 2026-05-19 to 2026-05-21 ($3$ consecutive full days,
$\approx 2.17$ million ticks per symbol). The robustness study (Section~\ref{sec:validation})
spans $31$ dates from 2026-04-19 to 2026-06-04, of which $23$ are sufficiently complete.

The feature vector spans fifteen categories computed directly from the book and trade stream:
raw book state, order-book imbalance, trade flow, realized and range volatility, permutation and
tick entropy, exchange context (funding, open interest), trend and Hurst statistics, illiquidity
(Kyle, Amihud, Roll), toxicity (VPIN, effective/realized spread), and derived composites. Of the
$236$ schema features, $207$ are populated in this study; the remainder depend on a
position-attribution pipeline outside our scope.

\subsection{The autonomous platform as apparatus}

Discovery is performed by an autonomous research system organized around a four-stage loop,
\[
  \text{Ingest (Rust)} \;\longrightarrow\; \text{Discover (agents)} \;\longrightarrow\;
  \text{Replicate (5 gates)} \;\longrightarrow\; \text{Deploy (paper)},
\]
with a feedback path that retires decayed signals. Four research agents operate at different
timescales (tick, minute-to-hour, hour-to-day, and a cross-agent orchestrator); each generates
hypotheses, tests them through the replication protocol of Section~\ref{sec:methodology}, and
registers only survivors. The design goal is not trading per se but \emph{disciplined discovery}:
the apparatus is engineered so that a feature reaching the registry has cleared independent
replication and FDR control, making the empirical claims below resistant to data-snooping.

\subsection{The information coefficient}

\begin{definition}[Forward return and information coefficient]
\label{def:ic}
Let $p_t$ be the mid-price and $h$ a forecast horizon. The forward return in basis points is
\begin{equation}
  r_{t,h} := \left(\frac{p_{t+h}}{p_t} - 1\right)\times 10^4 .
  \label{eq:fwd}
\end{equation}
The \emph{directional} information coefficient of a feature $f$ is the Spearman rank correlation
$\IC = \corr_{\mathrm{S}}(f_t,\, r_{t,h})$; the \emph{volatility} information coefficient replaces
$r_{t,h}$ by $|r_{t,h}|$.
\end{definition}

Spearman correlation is chosen for robustness to the heavy tails of tick returns. To suppress the
strong autocorrelation of $100$\,ms sampling, evaluation points are subsampled every $100$ ticks
($\approx 10$\,s), giving $\approx 21{,}700$ near-independent observations per symbol-horizon
cell; significance is reported at $p<0.01$ (two-sided), and significant cells are marked $*$.
Horizons examined are $\{1\text{s},5\text{s},30\text{s},1\text{m},5\text{m},15\text{m}\}$.

% ============================================================
\section{Disciplined Discovery Methodology}
\label{sec:methodology}
% ============================================================

\subsection{The five-gate replication protocol}

Every candidate signal must independently pass five gates before registration; failure at any
gate routes the hypothesis to a documented ``graveyard'' with its rejection reason. The protocol
is designed to control the false discovery rate under heavy multiple testing.

\begin{enumerate}
  \item[\textbf{G1}] \textbf{Discovery.} The rank IC must exceed an adaptive threshold and its
        incremental component $\mathrm{dIC}\geq 0.05$ (the marginal IC beyond already-registered
        signals).
  \item[\textbf{G2}] \textbf{Cost.} Gross return per trade must exceed a fee-based floor
        ($\geq 0.1$\,bps), so that statistically real signals too small to overcome costs are
        rejected.
  \item[\textbf{G3}] \textbf{Temporal replication.} The IC must hold on at least one additional
        date (a White, 2000, reality-check in spirit).
  \item[\textbf{G4}] \textbf{Cross-symbol replication.} The IC must hold on at least one other
        symbol.
  \item[\textbf{G5}] \textbf{Correlation deduplication.} The signal's correlation with every
        registered signal must satisfy $\rho < 0.7$, enforcing breadth.
\end{enumerate}

\subsection{FDR control, adaptive thresholds, and decay}

At the end of each discovery cycle, Benjamini--Hochberg FDR control (Benjamini and Hochberg, 1995)
is applied across all tests at $q=0.05$: order the $p$-values
$p_{(1)}\leq\cdots\leq p_{(m)}$, find $j^*=\max\{j: p_{(j)}\leq jq/m\}$, and retain the first
$j^*$. The acceptance threshold itself adapts as the registry strengthens,
\begin{equation}
  \mathrm{minIC}(t) = \max\!\Big(\,\text{floor},\;
  \med\{\IC_i : i\in R(t)\}\times 0.8\,\Big),
  \label{eq:adaptive}
\end{equation}
with a floor of $0.10$, so that a marginal signal cannot enter a registry already populated by
strong ones. Registered signals are continuously monitored and auto-retired if their rolling IC
falls below half of their discovery IC for fourteen consecutive cycles. Together these mechanisms
make the registry a conservative, self-pruning set rather than a catalogue of in-sample maxima.

\subsection{Information-theoretic feature selection}

As a nonlinear complement to rank correlation, an information-theoretic engine estimates mutual
information $I(f;r)$ via the Kraskov--St\"ogbauer--Grassberger $k$-nearest-neighbor estimator,
conditional mutual information $I(f;r\mid H)$, the interaction information
$\mathrm{II}(f;r;H)=I(f;r\mid H)-I(f;r)$ (positive for synergy, negative for redundancy), and a
linear transfer entropy. A feature is deemed cost-viable only if its mutual information clears a
fee-based minimum,
\begin{equation}
  I_{\min} = -\tfrac{1}{2}\log_2\!\Big(1 - (\text{fee}/\sigma_r)^2\Big),
  \label{eq:imin}
\end{equation}
the information required to overcome round-trip cost relative to return volatility $\sigma_r$.
Greedy forward selection then maximizes conditional MI gain, $f_{k+1}=\argmax_f I(f;r\mid S_k)$,
stopping when the marginal gain falls below $I_{\min}$---selecting features that are
uncorrelated-by-construction in the information metric.

% ============================================================
\section{Signal Discovery}
\label{sec:discovery}
% ============================================================

\subsection{Information-coefficient horizon profile}

Table~\ref{tab:horizon} reports directional IC versus horizon for the leading features on BTC.
Order-book imbalance (level-1 and depth-weighted) is the strongest directional predictor in the
entire feature set, peaking at $\IC\approx 0.45$ at $1$--$5$\,s; positive imbalance (more bids than
asks) predicts an up-move within seconds, and ask depth carries the mirror-image negative IC. The
half-life is $\approx 30$\,s: by five minutes roughly $80\%$ of the predictive content is gone. By
contrast, the spread and a slow regime-divergence feature carry IC that \emph{grows} with horizon,
capturing slow accumulation rather than microstructure---this weak residual is exactly what the
original $100$-minute evaluation was inadvertently measuring.

\begin{table}[ht]
\begin{center}
\small
\begin{tabular}{lrrrrrr}
\hline
\textbf{Feature} & \textbf{1s} & \textbf{5s} & \textbf{30s} & \textbf{1m} & \textbf{5m} & \textbf{15m} \\
\hline
\code{imbalance_qty_l1}        & $+0.447^{*}$ & $+0.453^{*}$ & $+0.257^{*}$ & $+0.184^{*}$ & $+0.089^{*}$ & $+0.056^{*}$ \\
\code{imbalance_depth_weighted}& $+0.427^{*}$ & $+0.434^{*}$ & $+0.255^{*}$ & $+0.185^{*}$ & $+0.081^{*}$ & $+0.037^{*}$ \\
\code{raw_ask_depth_5}         & $-0.414^{*}$ & $-0.425^{*}$ & $-0.238^{*}$ & $-0.167^{*}$ & $-0.072^{*}$ & $-0.040^{*}$ \\
\code{flow_aggressor_ratio_5s} & $+0.112^{*}$ & $+0.112^{*}$ & $+0.074^{*}$ & $+0.051^{*}$ & $+0.009$     & $-0.020$     \\
\code{raw_spread_bps}          & $-0.000$     & $+0.006$     & $+0.014$     & $+0.023^{*}$ & $+0.073^{*}$ & $+0.139^{*}$ \\
\code{regime_divergence_1h}    & $-0.020$     & $-0.009$     & $-0.035^{*}$ & $-0.043^{*}$ & $-0.058^{*}$ & $-0.074^{*}$ \\
\hline
\end{tabular}
\end{center}
\caption{Directional IC by horizon, BTC. Imbalance peaks at $1$--$5$\,s and decays; the spread and
regime divergence build slowly. $*: p<0.01$.}
\label{tab:horizon}
\end{table}

\subsection{Cross-symbol scan and signal universality}

Scanning all live features across BTC, ETH, and SOL, the order-book imbalance cluster dominates
directional prediction uniformly: peak IC of $0.40$--$0.47$ at $1$--$5$\,s on every symbol
(Table~\ref{tab:topdir}). The decay rate is symbol-dependent---half-lives of $\approx 30$\,s
(BTC), $\approx 20$\,s (ETH), and $\approx 15$\,s (SOL)---with the thinner book discovering price
faster. The hierarchy of features is identical across symbols, indicating a structural rather than
asset-specific phenomenon.

\begin{table}[ht]
\begin{center}
\small
\begin{tabular}{lrrrc}
\hline
\textbf{Feature} & \textbf{BTC} & \textbf{ETH} & \textbf{SOL} & \textbf{Horizon} \\
\hline
\code{imbalance_qty_l1}         & $+0.453$ & $+0.447$ & $+0.466$ & 1--5s \\
\code{imbalance_qty_l5}         & $+0.456$ & $+0.431$ & $+0.453$ & 5s \\
\code{imbalance_orders_l5}      & $+0.435$ & $+0.438$ & $+0.455$ & 1--5s \\
\code{imbalance_depth_weighted} & $+0.434$ & $+0.410$ & $+0.425$ & 5s \\
\code{raw_ask_depth_5}          & $-0.424$ & $-0.404$ & $-0.415$ & 5s \\
\code{raw_bid_depth_5}          & $+0.420$ & $+0.405$ & $+0.406$ & 5s \\
\code{cross_obi_mean}           & $+0.354$ & $+0.334$ & $+0.331$ & 1--5s \\
\code{micro_queue_position_bid} & $+0.307$ & $+0.282$ & $+0.291$ & 5s \\
\code{flow_vwap_deviation}      & $-0.287$ & $-0.206$ & $-0.188$ & 1s \\
\code{micro_obi_velocity}       & $+0.192$ & $+0.181$ & $+0.208$ & 1s \\
\code{flow_aggressor_ratio_5s}  & $+0.112$ & $+0.109$ & $+0.111$ & 1--5s \\
\hline
\end{tabular}
\end{center}
\caption{Top directional features at their $1$--$5$\,s peak, all $p<0.01$. The imbalance cluster
dominates across all three symbols.}
\label{tab:topdir}
\end{table}

\subsection{Eight independent directional axes}

Within the highly-correlated directional cluster, the scan resolves eight approximately
independent axes (Table~\ref{tab:axes}); these are the inputs combined in
Section~\ref{sec:combiner}.

\begin{table}[ht]
\begin{center}
\small
\begin{tabular}{llr}
\hline
\textbf{Axis} & \textbf{Representative feature} & \textbf{IC ($1$--$5$s)} \\
\hline
Order-book imbalance       & \code{imbalance_qty_l1}             & $0.45$ \\
Raw depth asymmetry        & \code{raw_bid_depth_5}              & $0.42$ \\
Cross-symbol imbalance     & \code{cross_obi_mean}               & $0.35$ \\
Queue dynamics             & \code{micro_queue_position_bid}     & $0.31$ \\
VWAP deviation (reversion) & \code{flow_vwap_deviation}          & $0.25$ \\
OBI velocity (momentum)    & \code{micro_obi_velocity}           & $0.19$ \\
Imbalance entropy          & \code{ent_permutation_imbalance_16} & $0.17$ \\
Aggressor flow             & \code{flow_aggressor_ratio_5s}      & $0.11$ \\
\hline
\end{tabular}
\end{center}
\caption{Eight approximately independent directional axes and representative features.}
\label{tab:axes}
\end{table}

\subsection{Directional/volatility orthogonality}

Direction and magnitude are measured by disjoint feature sets. Imbalance features have
essentially zero volatility IC; the best volatility predictors---self-exciting intensity
(\code{hawkes_intensity}, $\IC\approx 0.35$ at $5$\,s), realized volatility, trade intensity, and
effective spread---have essentially zero directional IC. This orthogonality is consistent across
all three symbols and motivates the separation of \emph{direction}, \emph{timing}, and
\emph{sizing} in the combiner of Section~\ref{sec:combiner}.

% ============================================================
\section{Robustness Validation}
\label{sec:validation}
% ============================================================

We subject the candidate signal---order-book imbalance at $1$--$5$\,s---to eight checks. Seven are
confirmatory; the eighth (the conditional IC of Section~\ref{sec:adverse}) is the decisive
constraint.

\paragraph{Stability across days, hours, and regimes.}
Per-day IC is stable across $23$ valid dates (Table~\ref{tab:perday}). The signal is present
$24$/$7$ with no dead intraday zone (a slight dip at 12--16 UTC, strongest at 04--08 UTC).
Splitting on median five-minute realized volatility, IC remains above $0.34$ in both regimes for
all symbols, and is slightly stronger in low volatility---imbalance is more informative when the
book is not being disrupted.

\begin{table}[ht]
\begin{center}
\small
\begin{tabular}{lrrr}
\hline
\textbf{Horizon} & \textbf{BTC} & \textbf{ETH} & \textbf{SOL} \\
\hline
1s  & $+0.442 \pm 0.069$ & $+0.418 \pm 0.099$ & $+0.415 \pm 0.052$ \\
5s  & $+0.433 \pm 0.061$ & $+0.400 \pm 0.098$ & $+0.355 \pm 0.071$ \\
30s & $+0.260 \pm 0.068$ & $+0.224 \pm 0.072$ & $+0.182 \pm 0.048$ \\
\hline
\end{tabular}
\end{center}
\caption{Per-day rolling IC of \code{imbalance_qty_l1} (mean $\pm$ standard deviation across days).}
\label{tab:perday}
\end{table}

\paragraph{Bootstrap and temporal out-of-sample.}
A $1000$-resample bootstrap on the discovery window gives $95\%$ confidence intervals of width
$\approx 0.02$: for \code{imbalance_qty_l1} at $5$\,s, BTC $[0.442,0.465]$, ETH $[0.425,0.450]$,
SOL $[0.400,0.425]$. The one cautionary check is temporal out-of-sample
(Table~\ref{tab:oos}): comparing the early window (May 19--21) to the late window (Jun 2--4), the
$5$\,s-horizon IC drops $\approx 0.15$ on every symbol while the $1$\,s-horizon IC is stable or
improves (BTC $0.447\to 0.502$). June dates carry wider spreads, consistent with thinner books and
faster price discovery: the signal is not disappearing, its half-life is shortening.

\begin{table}[ht]
\begin{center}
\small
\begin{tabular}{lcccccc}
\hline
\textbf{Feature ($5$s)} & \textbf{BTC early} & \textbf{BTC late} & \textbf{ETH early} & \textbf{ETH late} & \textbf{SOL early} & \textbf{SOL late} \\
\hline
\code{imbalance_qty_l1}         & $+0.453$ & $+0.286$ & $+0.438$ & $+0.277$ & $+0.412$ & $+0.240$ \\
\code{imbalance_depth_weighted} & $+0.435$ & $+0.317$ & $+0.395$ & $+0.270$ & $+0.377$ & $+0.242$ \\
\code{raw_ask_depth_5}          & $-0.425$ & $-0.282$ & $-0.395$ & $-0.261$ & $-0.367$ & $-0.233$ \\
\hline
\end{tabular}
\end{center}
\caption{Temporal out-of-sample IC at $5$\,s (early $\to$ late). The $5$\,s IC weakens while the
$1$\,s IC holds---a shortening half-life.}
\label{tab:oos}
\end{table}

\paragraph{Validation matrix.}
Aggregating the checks against fixed \emph{a priori} thresholds (Table~\ref{tab:matrix}), the
signal passes seven of eight for BTC and six of eight for ETH/SOL. The failures are the
temporal-OOS decay and, for ETH/SOL, the worst-single-day floor.

\begin{table}[ht]
\begin{center}
\small
\begin{tabular}{lcccl}
\hline
\textbf{Check} & \textbf{BTC} & \textbf{ETH} & \textbf{SOL} & \textbf{Threshold} \\
\hline
Per-day IC mean ($5$s)   & $0.43$ P & $0.40$ P & $0.35$ P & $>0.30$ \\
Per-day IC std ($5$s)    & $0.06$ P & $0.10$ P & $0.07$ P & $<0.10$ \\
Worst single day ($5$s)  & $0.24$ P & $0.04$ F & $0.17$ F & $>0.20$ \\
Intraday (all windows)   & P & P & P & all $>0.20$ \\
Low-volatility IC        & $0.46$ P & $0.43$ P & $0.39$ P & $>0.30$ \\
High-volatility IC       & $0.42$ P & $0.39$ P & $0.34$ P & $>0.30$ \\
Bootstrap CI lower       & $0.44$ P & $0.43$ P & $0.40$ P & $>0.30$ \\
Temporal OOS delta       & $0.17$ F & $0.16$ F & $0.17$ F & $<0.10$ \\
\hline
\textbf{Total}           & \textbf{7/8} & \textbf{6/8} & \textbf{6/8} & \\
\hline
\end{tabular}
\end{center}
\caption{Validation matrix. The signal is real, stable, and universal; the only systematic failure
is the shortening horizon.}
\label{tab:matrix}
\end{table}

% ============================================================
\section{The Adverse-Selection Barrier}
\label{sec:adverse}
% ============================================================

The seven confirmatory checks establish a strong, robust, universal predictor. The eighth check
asks the economic question: can a passive liquidity provider capture it?

\begin{definition}[Conditional information coefficient]
\label{def:cond}
Let $\mathcal{F}$ be the event that a directionally-correct passive limit order---a buy at the best
bid when imbalance predicts up, or a sell at the best ask when it predicts down---fills within a
fixed tick budget. The \emph{conditional} information coefficient is the IC computed on the
subsample $\{t : \mathcal{F} \text{ occurs}\}$.
\end{definition}

Table~\ref{tab:cond} reports the result. The unconditional IC is $0.41$--$0.45$. Conditioning on
the \emph{any}-fill event \emph{raises} it (to $0.46$--$0.53$), because states near fills carry
larger absolute imbalance and hence larger returns. But conditioning on the
\emph{directionally-correct} fill---a buy that actually fills when imbalance predicts up---collapses
the IC to $\approx 0.03$ on every symbol.

\begin{table}[ht]
\begin{center}
\small
\begin{tabular}{lrrr}
\hline
\textbf{Condition} & \textbf{BTC} & \textbf{ETH} & \textbf{SOL} \\
\hline
Unconditional                    & $+0.453$ & $+0.438$ & $+0.412$ \\
Any fill event                   & $+0.526$ & $+0.516$ & $+0.460$ \\
Buy fill ($\text{imb}>0$)        & $+0.032$ & $-0.061$ & $-0.032$ \\
Sell fill ($\text{imb}<0$)       & $-0.047$ & $-0.027$ & $-0.021$ \\
\hline
\end{tabular}
\end{center}
\caption{Conditional IC of \code{imbalance_qty_l1}. The directionally-correct passive fill destroys
the edge.}
\label{tab:cond}
\end{table}

\begin{proposition}[Structural adverse selection]
\label{prop:adverse}
Under a passive mid-cross fill model, a directional order-book-imbalance signal is not monetizable:
the conditional IC on directionally-correct fills is statistically indistinguishable from zero.
\end{proposition}

\noindent The mechanism is structural, not statistical:
\begin{enumerate}
  \item The signal says ``buy'' (imbalance $>0$); a passive buy is posted at the best bid.
  \item The fill requires the mid to fall to the bid---an adverse move.
  \item At the instant of the fill, the ``up'' prediction has already been invalidated; the
        post-fill edge is zero.
\end{enumerate}
Fill prevalence rises with book thinness: directionally-correct buy/sell fills occur on
$20.8\%/21.7\%$ of BTC states, $24.1\%/24.6\%$ of ETH states, and $30.8\%/31.1\%$ of SOL states.
This explains the otherwise paradoxical observation that a naive passive simulator generates
\emph{more} trades on SOL but \emph{worse} performance---more fills means more adverse selection.
Proposition~\ref{prop:adverse} is the empirical content of Glosten--Milgrom (1985) adverse
selection, measured directly at the per-fill level.

% ============================================================
\section{Hierarchical Signal Combination}
\label{sec:combiner}
% ============================================================

If passive timing alone is non-monetizable, the natural response is to trade only where the
expected move exceeds execution cost. We separate the eleven leading features into three
functional layers and compose them multiplicatively.

\begin{definition}[Three-layer composite]
\label{def:combiner}
Let $L_1\in\{-1,0,+1\}$ be a slow directional bias (from regime divergence, spread, and a
short trend EMA), active only when its IC-weighted $z$-score exceeds a band; let $L_2\in[-1,1]$ be
a fast entry signal (from imbalance, cross-symbol imbalance, queue position, VWAP deviation, and
OBI velocity), \emph{zeroed when its sign disagrees with $L_1$}; and let $L_3\in[0,1]$ be an
inverse-volatility scale (from self-exciting intensity, trade count, and realized volatility). The
composite is
\begin{equation}
  S = L_1 \cdot |L_2| \cdot L_3,
  \qquad
  S = 0.3\,|L_2|\,L_3 \text{ when } L_1 \text{ is neutral.}
  \label{eq:composite}
\end{equation}
\end{definition}

The directional gate is the key device: fast signals alone are non-monetizable (conditional IC
$\approx 0.03$), but fast signals \emph{confirming} a slow bias operate at horizons where the
expected move exceeds cost; inverse-volatility sizing reduces exposure precisely when adverse
selection on passive orders is worst. Trained with a four-fold expanding window and a $100$-bar
embargo on five-minute bars, the combiner attains positive cost-adjusted out-of-sample
performance (Table~\ref{tab:combiner}); the learned IC weights track the in-scan priors, arguing
against pure walk-forward overfitting.

\begin{table}[ht]
\begin{center}
\small
\begin{tabular}{lrrr}
\hline
\textbf{Metric (mean over folds)} & \textbf{BTC} & \textbf{ETH} & \textbf{SOL} \\
\hline
Composite IC          & $+0.178$ & $+0.248$ & $+0.359$ \\
Directional accuracy  & $0.557$  & $0.576$  & $0.594$  \\
Sharpe (cost-adjusted)& $+1.25$  & $+1.71$  & $+2.40$  \\
\hline
\end{tabular}
\end{center}
\caption{Hierarchical combiner walk-forward out-of-sample results (BTC/ETH at a $5$h horizon, SOL
at $3.3$h owing to less data).}
\label{tab:combiner}
\end{table}

\paragraph{Honest assessment.}
These results rest on a thin two-day window, so the confidence intervals are wide; the monotone
improvement across folds is a caution flag for look-ahead or a trending sample; the slow layer
dominates, so the marginal value of the fast layers over the slow bias alone is not yet proven;
and the SOL figures (shorter horizon, less data) are an upper bound. The result is presented as a
proof of concept that adverse selection can be \emph{designed around}, not as a deployable
strategy. The required next steps are data accumulation, a slow-layer-only ablation, and
conditional-fill analysis on real execution logs.

% ============================================================
\section{Discussion}
\label{sec:discussion}
% ============================================================

The results invert the usual order of difficulty. Discovering a strong, robust, universal
directional predictor proved straightforward; the entire economic question reduces to whether
that predictor survives adverse selection. Three implications follow.

\paragraph{Execution is the frontier.} The conditional-IC collapse (Section~\ref{sec:adverse}) is a
property of the \emph{fill model}, not the signal. A passive mid-cross fill is the worst case; an
execution model built from actual aggressor trades may recover part of the edge, and measuring the
conditional IC under such a model---rather than the unconditional IC---is the correct objective
function for microstructure alpha research. Statistical significance of a predictor is necessary
but not sufficient; the relevant quantity is the IC \emph{conditional on the states one can
transact in}.

\paragraph{The signal is perishable.} The $\approx 0.15$ five-second decay from May to June, with a
stable one-second IC, indicates that the tradeable horizon is contracting as books thin. Any
deployment must track this and adapt per symbol, with SOL fastest (half-life $\approx 15$\,s).

\paragraph{Methodology as contribution.} The empirical claims are credible precisely because they
are the output of a false-discovery-controlled apparatus: five independent replication gates, FDR
control, an adaptive threshold, and continuous decay monitoring. We suggest this design---automated
discovery with replication and FDR built into the loop rather than applied after the fact---as a
template for empirical microstructure research, which is otherwise highly exposed to the
multiple-testing critique of Harvey, Liu, and Zhu (2016).

% ============================================================
\section{Conclusion}
\label{sec:conclusion}
% ============================================================

Order-book imbalance predicts the direction of Hyperliquid perpetual returns with
$\IC\approx 0.45$ at $1$--$5$ seconds, universally across BTC, ETH, and SOL, surviving an
eight-check robustness battery administered by an autonomous, false-discovery-controlled discovery
platform. The same data show, decisively, that a passive mid-cross fill model eliminates the edge
through adverse selection: the conditional IC collapses to $\approx 0.03$. A hierarchical combiner
that gates fast timing on a slow directional bias recovers positive cost-adjusted out-of-sample
performance on a thin window, demonstrating that adverse selection can be confronted by design. The
binding constraint of microstructure alpha is therefore not signal quality but execution
feasibility, and the right object of study is the information coefficient conditional on
transactable states. Future work will replace the passive fill model with an event-driven model
built from per-trade data, extend the robustness study as data accumulate, and test the combiner's
slow-layer ablation.

% ============================================================
\section*{Reproducibility}
\addcontentsline{toc}{section}{Reproducibility}
% ============================================================

The data are Hyperliquid Level-2 order-book and trade streams for BTC, ETH, and SOL at $100$\,ms,
ingested in Rust to columnar storage. The discovery window is 2026-05-19 to 2026-05-21; the
robustness window is 2026-04-19 to 2026-06-04 ($23$ valid dates). The information coefficient is
the Spearman rank correlation of Definition~\ref{def:ic}, subsampled every $100$ ticks, with
significance at $p<0.01$. Five-gate replication, Benjamini--Hochberg FDR at $q=0.05$, and the
adaptive threshold of equation~\eqref{eq:adaptive} are applied as described in
Section~\ref{sec:methodology}. Exact code revision (git SHA), per-run data fingerprints, and the
feature manifest are recorded with each result and available on request. No trading thresholds,
entry logic, or position-sizing parameters are reported; the paper concerns signal quality and its
conditional behavior, not a deployable strategy.

% ============================================================
\section*{References}
\addcontentsline{toc}{section}{References}
% ============================================================

\begin{enumerate}
  \item A\"it-Sahalia, Y.\ \& Brunetti, C.\ (2024). What do we know about cryptocurrency?
        \textit{Annual Review of Financial Economics}, 16.

  \item Bailey, D.H.\ \& L\'opez de Prado, M.\ (2014). The deflated Sharpe ratio: correcting for
        selection bias, backtest overfitting, and non-normality. \textit{Journal of Portfolio
        Management}, 40(5), 94--107.

  \item Benjamini, Y.\ \& Hochberg, Y.\ (1995). Controlling the false discovery rate: a practical
        and powerful approach to multiple testing. \textit{Journal of the Royal Statistical
        Society, Series B}, 57(1), 289--300.

  \item Cont, R., Kukanov, A.\ \& Stoikov, S.\ (2014). The price impact of order book events.
        \textit{Journal of Financial Econometrics}, 12(1), 47--88.

  \item Cont, R., Stoikov, S.\ \& Talreja, R.\ (2010). A stochastic model for order book dynamics.
        \textit{Operations Research}, 58(3), 549--563.

  \item Easley, D., L\'opez de Prado, M.M.\ \& O'Hara, M.\ (2012). Flow toxicity and liquidity in a
        high-frequency world. \textit{Review of Financial Studies}, 25(5), 1457--1493.

  \item Glosten, L.R.\ \& Milgrom, P.R.\ (1985). Bid, ask and transaction prices in a specialist
        market with heterogeneously informed traders. \textit{Journal of Financial Economics},
        14(1), 71--100.

  \item Harvey, C.R., Liu, Y.\ \& Zhu, H.\ (2016). $\ldots$ and the cross-section of expected
        returns. \textit{Review of Financial Studies}, 29(1), 5--68.

  \item Kraskov, A., St\"ogbauer, H.\ \& Grassberger, P.\ (2004). Estimating mutual information.
        \textit{Physical Review E}, 69(6), 066138.

  \item Kyle, A.S.\ (1985). Continuous auctions and insider trading. \textit{Econometrica}, 53(6),
        1315--1335.

  \item Lee, S.S.\ \& Mykland, P.A.\ (2008). Jumps in financial markets: a new nonparametric test
        and jump dynamics. \textit{Review of Financial Studies}, 21(6), 2535--2563.

  \item Makarov, I.\ \& Schoar, A.\ (2020). Trading and arbitrage in cryptocurrency markets.
        \textit{Journal of Financial Economics}, 135(2), 293--319.

  \item Schreiber, T.\ (2000). Measuring information transfer. \textit{Physical Review Letters},
        85(2), 461--464.

  \item White, H.\ (2000). A reality check for data snooping. \textit{Econometrica}, 68(5),
        1097--1126.
\end{enumerate}

\end{document}
